\section{Experimental design in full}\label{sec:design}

The comparison set, the primary metrics, the cross-validation protocol and the families of statistical tests were fixed in a written study plan before the benchmark was run (the plan is in \texttt{docs/study\_plan.md} of the repository). The width of the region of practical equivalence, several of the mechanism and dataset-level analyses and all additional experiments and analyses (Section~\ref{sec:design:revision}), which were added after the first results were available, are exploratory; Table~\ref{tab:analysis_status} lists the status of every analysis.

\subsection{Protocol}\label{sec:design:protocol}

Each method was evaluated by five-times repeated stratified five-fold cross-validation (25 train--test splits per dataset), with fixed seeds and identical splits for all methods. Within each training fold, median imputation and min--max scaling to $[0,1]$ were fitted on the training data and applied to the test fold; the sampler was then applied to the scaled training fold, and the classifier was trained on the resampled data and evaluated on the untouched test fold. No information from a test fold entered resampling or training.

\subsection{Compared methods}\label{sec:design:methods}

The comparison set contains sixteen undersampling-family configurations: no resampling (\textsc{none}); eleven standard undersamplers from imbalanced-learn~0.14.2 \citep{lemaitre2017imbalanced} with their default settings---random undersampling (\RUS), NearMiss versions 1, 2 and 3 (3 neighbours), Tomek links, edited nearest neighbours (3 neighbours), the neighbourhood cleaning rule, one-sided selection, condensed nearest neighbour, the instance hardness threshold (random forest with 100 trees, 5-fold cross-validation) and ClusterCentroids with $k$-means and hard voting, so that real instances are retained (a deviation from the library default, which uses soft voting, i.e. synthetic centroids, for dense data); and the four matching undersamplers. The matching undersamplers target a balanced training set ($m=p$), \MUSCS{} and \MUSCO{} keep a random member of each pair, and \MUSCO{} uses $k=10$. Random undersampling, NearMiss, ClusterCentroids and the matching undersamplers thus all target a balanced training set; we call them the \emph{balancing} methods. All of them reach the target exactly except NearMiss-3, whose first stage short-lists the three nearest majority neighbours of each minority instance and can therefore return fewer instances: on 26 of the 55 datasets it retained less than the target (median 76\% of it on these datasets, as little as 12\%). The cleaning methods do not control the output size and retain most majority instances (median 16 times the minority size), one-sided selection about 8 times, and condensed nearest neighbour a consistent subset of very variable size (median 0.97 times the minority size, range 0.02 to 6.6). The instance hardness threshold nominally targets a balanced output, but because it keeps every majority instance whose out-of-fold probability reaches a percentile threshold, ties at that threshold make it keep a median of 12 times the minority size. Two non-undersampling references are reported separately: SMOTE with five neighbours \citep{chawla2002smote} and cost-sensitive learning with class weights inversely proportional to class frequencies.

\subsection{Classifiers and metrics}\label{sec:design:classifiers}

We use five classifiers from scikit-learn \citep{pedregosa2011scikit} with default settings: CART decision trees \citep{breiman2017classification}, $k$-nearest neighbours with $k=5$ \citep{cover1967nearest}, support vector machines with an RBF kernel, $C=1$ and \texttt{gamma=scale} \citep{cortes1995support}, $L_2$-regularised logistic regression with $C=1$, and random forests with 100 trees \citep{breiman2001random}. Hyperparameters were not tuned, so that differences reflect the training data rather than tuning effort; this choice is discussed in Section~6 of the main paper. The two primary metrics are the geometric mean of the class-wise recalls, $\text{G-mean}=\sqrt{\mathrm{TPR}\cdot\mathrm{TNR}}$ \citep{kubat1997addressing,kubat1998machine}, which evaluates the default decision threshold, and the area under the ROC curve, which evaluates the ranking of test instances. Balanced accuracy, the F1 score of the minority class, the Matthews correlation coefficient \citep{chicco2020advantages} and the area under the precision--recall curve \citep{saito2015precision} are secondary metrics. Scores for the ROC and precision--recall curves are predicted probabilities, or the decision function for the support vector machine. Two caveats apply. G-mean at the default threshold mixes the ranking quality of a classifier with the placement of its scores relative to that threshold, so Section~\ref{sec:design:revision} adds a threshold-adjusted analysis. And ROC-AUC is threshold-free only for classifiers with fine-grained scores: unpruned CART trees output leaf frequencies that are mostly 0 or 1, so that their ROC-AUC is close to the balanced accuracy at the default threshold, and 5-NN produces at most six distinct scores. With 2 to 4 minority instances per test fold on the smallest datasets, fold-level metrics are noisy; as a sensitivity analysis we also compute G-mean, F1 and MCC from confusion matrices pooled over the five folds of each repetition \citep{forman2010apples}.

\subsection{Statistical analysis}\label{sec:design:stats}

For each dataset, method and classifier, fold-level results were averaged over the 25 splits. For each classifier and primary metric we report average ranks over the 55 datasets, the Friedman test with the Iman--Davenport correction \citep{friedman1937use,iman1980approximations,demsar2006statistical} and critical-difference diagrams based on the Nemenyi test at $\alpha=0.05$. Because mean-rank post-hoc tests depend on the pool of compared methods \citep{benavoli2016should}, inferential pairwise conclusions are based on two-sided Wilcoxon signed-rank tests \citep{wilcoxon1945individual} of each matching undersampler against each of the twelve non-matching configurations, with Holm's correction \citep{holm1979simple} over the 48 comparisons per classifier and metric. The 2$\times$2 contrasts (stable versus optimal within each arm, boundary versus coverage within each solution concept) are tested in the same way. We designate the 40 comparisons of the four matching undersamplers with \RUS{} for the five classifiers and the two primary metrics as the primary family and also report Holm's correction over these 40 tests; because this family was designated after the results were known, it is not confirmatory in the strict sense. All other tests are exploratory, and their Holm corrections control the error rate only within the stated family. We complement these tests with Bayesian signed-rank tests \citep{benavoli2017time} with a region of practical equivalence (ROPE) of $\pm 0.01$ on the metric scale, the value used for accuracy by \citet{benavoli2017time}; because this width is a convention rather than a pre-specified, application-based choice, we also report the posteriors for widths of 0.0025, 0.005 and 0.02 and relate the width to the cross-validation noise of a single dataset, estimated by the corrected resampled-$t$ standard error \citep{nadeau2003inference}. Besides random undersampling and the best non-matching method, the Bayesian tests also compare the matching undersamplers with ClusterCentroids. The best non-matching method of each classifier is selected by its average rank on the same data that are used to test the matching undersamplers against it, which biases these comparisons against the matching undersamplers. Datasets of the same family share base data and are not independent; besides the one-per-family subset below, we compute family-clustered bootstrap intervals (24 families resampled with replacement) for the main paired differences. Robustness is assessed on (a) a pre-specified subset of 24 datasets containing one dataset per family (the member with median imbalance ratio), (b) the benchmark without the nine datasets that are perfectly separable by a single feature, (c) the low- and high-IR groups separately, and (d) the secondary metrics.

\subsection{Ablations}\label{sec:design:ablations}

With CART, $k$-NN and SVM, we varied (i) the target ratio $m/p\in\{1,1.5,2\}$ for \RUS, NearMiss-1, ClusterCentroids, \MUSBS, \MUSBO{} and \MUSCS{} (targets are capped at the available majority size; \MUSCO{} was not ablated because of its cost), (ii) the keep rule of \MUSCS{} (random versus closer to the minority class) and (iii) the preference structure of \MUSBS{} (Euclidean versus class-conditional Mahalanobis). The effect of the candidate degree $k$ of \MUSCO{} was assessed directly through the matching cost gap to the complete-graph optimum.

\subsection{Additional experiments}\label{sec:design:revision}

Six further experiments test whether the main conclusions depend on choices of the protocol. They use the eight balancing methods other than \MUSCO, whose cost (Section~\ref{sec:results:runtime}) made repeated runs impractical; all other settings are those of the main experiment.
(i) \emph{Threshold adjustment and tuning}, run on the first three repetitions of the cross-validation (15 training folds per dataset) because of its cost. In every training fold, a stratified three-fold inner cross-validation, with the sampler applied inside each inner training part, provides out-of-fold scores. The threshold-adjusted variant keeps the default hyperparameters and uses the threshold that maximises G-mean on these scores; the tuned variant selects the hyperparameters with the best inner G-mean at the default threshold from small grids (CART: minimum leaf size 1, 2, 5 or 10; $k$-NN: $k\in\{1,3,5,9,15\}$; SVM: $C\in\{0.1,1,10,100\}$ and $\gamma$ equal to 0.1, 1 or 10 times the default; logistic regression: $C\in\{0.01,0.1,1,10,100\}$; random forests: minimum leaf size 1 or 3 and 50\% or the square root of the features per split).
(ii) \emph{Row order.} \MUSBS, \MUSBO{} and \MUSCS{} are rerun on a random permutation of the training rows, which changes the tie-breaking order; the selected instances are then sorted back into the original order, so that the classifiers see them in the same order as in the main experiment.
(iii) \emph{Intervention on the boundary arm.} Two variants of \MUSBS{} test the explanation of its sensitivity loss offered in Section~\ref{sec:results:mechanisms}: one gives capacity only to safe and borderline minority instances (\emph{safecap}; rare and outlier minority instances keep no majority partner but remain in the training set), the other declares majority instances with at least three minority instances among their five nearest neighbours unacceptable as partners (\emph{nounsafe}).
(iv) \emph{Spatially balanced baselines.} The local pivotal method in its LPM2 form \citep{grafstrom2012spatially}, applied to the majority class with equal inclusion probabilities and Euclidean distances on the preprocessed features, which returns exactly the target number of instances; and random undersampling stratified by a $k$-means partition of the majority class into ten clusters (fewer if the target is smaller) with proportional, largest-remainder allocation.
(v) \emph{Ensembles.} Ten balanced subsets per training fold, drawn by \RUS, \MUSCS{} or the local pivotal method with different seeds, each train one model, and the ensemble averages their scores; the ten single-subset models also give the variability of a single subset across seeds.
(vi) \emph{Duplicates.} On the datasets with exact duplicate rows (identical features and label), the duplicates are removed, new folds are drawn with the same seeds, and the balancing methods and no resampling are rerun.
In addition, we computed the realised Wasserstein distortion $\Wone(\mu_N,\mu_S)$ of all balancing methods and the coverage bound of Corollary~1 of the main paper on the 275 training folds of the first repetition, the ratio of the greedy to the optimal matching cost in the rounds of the coverage arm, the gap between the candidate graph $G_{10}$ and the complete graph on the benchmark folds with at most 1200 majority instances, family-clustered bootstrap intervals, metrics from pooled confusion matrices, Bayesian tests for four ROPE widths and against ClusterCentroids, and per-dataset differences with the corrected resampled-$t$ standard error \citep{nadeau2003inference}.

\begin{table}[t]
\centering
\caption{Status of the analyses. Planned: fixed in the study plan before the benchmark was run. Exploratory: chosen or added after the first results were available.}\label{tab:analysis_status}
\footnotesize
\begin{tabular}{@{}p{0.62\linewidth}l@{}}
\toprule
Analysis & Status \\
\midrule
Comparison set, classifiers, 5$\times$5 cross-validation, primary and secondary metrics & planned \\
Average ranks, Friedman/Iman--Davenport, Nemenyi diagrams; Wilcoxon--Holm comparisons with the non-matching methods; 2$\times$2 contrasts & planned \\
Bayesian signed-rank tests (ROPE width $\pm0.01$ chosen later) & planned (width exploratory) \\
Robustness: one dataset per family, IR groups, secondary metrics; ablations of target ratio, keep rule and preferences & planned \\
Removal of the nine single-feature-separable datasets & exploratory \\
Selection diagnostics: boundary proximity, distances to the minority and to the full majority, retention of unsafe majority instances & planned (definitions exploratory) \\
Coverage error, decomposition into sensitivity and specificity, coverage saturation & exploratory \\
Correlations of dataset-level gains with size, dimensionality, IR and complexity measures & planned \\
Safe-share threshold 0.8, depth-two regression tree & exploratory \\
Primary family of 40 comparisons with \RUS{} (designated after the first results) & exploratory \\
All analyses of Section~\ref{sec:design:revision} & exploratory \\
\bottomrule
\end{tabular}
\end{table}

\subsection{Implementation}\label{sec:design:implementation}

All experiments were run in Python~3.11 with NumPy~2.4.6 \citep{harris2020array}, SciPy~1.17.1 \citep{virtanen2020scipy}, scikit-learn~1.9.1, imbalanced-learn~0.14.2, NetworkX~3.6.1, pandas~3.1.1, \texttt{matching}~1.4.3 (used as a reference solver in the tests; \citealp{wilde2020matching}), scikit-posthocs~0.18.0 \citep{terpilowski2019scikit}, baycomp~1.0.2 (Bayesian signed-rank tests), statsmodels~0.15.0 (LOWESS), problexity~0.5.11 (complexity measures) and the dataset loader imbalanced-databases~0.1.1, on a computer with an Intel Core i9-13900HX processor (32 logical processors) and 32\,GB of memory running Windows, with independent processes per batch of train--test splits. The implementation of the four matching undersamplers passes 289 unit tests, which check stability, uniqueness and invariance (Propositions~1, 2(a) and~4 of the main paper), the ball characterisation, Wasserstein optimality against brute force, agreement with reference solvers and exact output sizes.
